\section{Research design}

We ask whether social media micro-influencers can be recruited to increase awareness and knowledge of misinformation among their online audiences. We investigate this question by partnering with Africa Check,\footnote{We conducted a small pilot with Chequea Bolivia in Bolivia.} who has recruited social media influencers to disseminate digital literacy training materials and fact-checks to their followers.

Small-scale social media influencers in Kenya and South Africa were first identified using the academic Twitter API, which allowed  researchers to collect information on all public Twitter accounts. These influencers were screened according to three inclusion criteria. The first criterion is location. We track all tweets based in Kenya and South Africa for approximately a year.\footnote{Twitter tags tweets with a \textit{Twitter Place}.} We then collected diverse information from the profiles of users, including location. While the Twitter API offers a location field associated with the user's profile, it has a free-form value, and thus may indicate a non-valid location or no location at all. Therefore, to identify social media users based in Kenya and South Africa, we filtered handles that explicitly referenced one of those countries as the user's location. As a result, we excluded profiles that do not have a study country as their location. The second criterion is the popularity of Twitter users. We select users with more than the country's average number of followers among our sample of Twitter users in Kenya and South Africa.\footnote{Kenya's average number of followers was 4,490, while South Africa's average number of followers was 3,265.}  The third criterion is whether the individual engages in high-quality news sharing. To this end, we retrieved the information shared by the users during a fixed period. Using a list of reputable media outlets in each country, which we validated with corresponding fact-checkers, we identify posts that contain links to a high-quality news article or fact-check. We only kept those users who had made at least one of such posts. 

In practice, our lists of social media micro-influencers comprise relatively journalists and social activists in each country. A natural concern with this pool of influencers is the possibility of floor effects, given that these influencers usually share high-quality information and their followers probably have higher social media literacy and are already less susceptible to misinformation. While this is a reasonable concern, our partnering fact-checking institutions are unwilling to recruit potential misinformers to disseminate the intervention content. Nevertheless, it remains unclear whether even the followers of more credible influencers yet possess sufficient social media literacy, while they are also the most likely to take up the intervention information. Applying the inclusion criteria yielded an initial sample of 732 from Kenya, and 1,309 from South Africa. The fact-checkers then vetted all proposed potential user media followers, filtering those they did not think would be good ambassadors for their institutions, ending up with a sample of close to 150 influencers across both countries. Finally, the fact-checkers approached these potential influencers to gauge interest in participating in the study.

Since this procedure did not yield sufficient social media micro-influencers interested in the project, it was complemented by relaxing the second and third inclusion criteria for the first batch. Relaxing the inclusion criteria yielded an additional sample of 1,890 influencers from Kenya, and 8,948 from South Africa. Twitter users who were recruited through these processes were vetted by Africa Check using similar criteria, ending up with a sample of 445 influencers from Kenya and  476 from South Africa. Additionally, for the second batch, we complemented our sample through social-media recruitment performed by Africa Check using Twitter ads.\footnote{A similar experience happened in Bolivia with Chequea Bolivia.} 

We have recruited a total of 302 social media influencers across both countries in what we denote three phases: \textbf{pilot} (40), \textbf{first batch} (152), and \textbf{second batch} (110).\footnote{Out of the 40 influencers recruited for the pilot, 20 influencers corresponded to each country. Out of the 152 influencers recruited for the first batch, 76 influencers corresponded to each country. Out of the 110 influencers recruited for the second batch, 58 influencers corresponded to Kenya and 52 corresponded to South Africa.} Upon enrollment into the study, we \textcolor{red}{block} randomized half of the social influencers within each context to receive the \textit{treatment} condition, while the other half served as the \textit{control} group\footnote{Refer to Appendix \ref{app:inf_rand} for more detail on this.}. Treated influencers received videos---which are produced by the relevant institutional partner---providing media literacy information. Treated influencers were further  encouraged to share one fact-check produced by the country's partnering fact-checker with their followers on social media. These prebunking and debunking efforts are designed to be complements, and have been disseminated by treated influencers for a period of eight weeks. 

\textcolor{red}{\bf Thanks for the response but this is an overkill for the paper. I can summarize this here, but can we have an Appendix table with this information?}
 All of the 40 influencers from the pilot stage were obtained from the sample generated using all three inclusion criteria. For the first batch, 61 influencers (47 from Kenya and 14 from South Africa) were obtained from the sample that took into account all the inclusion criteria, while 45 influencers (11 from Kenya and 34 from South Africa) were obtained by relaxing the third criterion, meaning that these users didn't share any high-quality news but had more followers than their country's average. Finally, 46 influencers (18 from Kenya and 28 from South Africa) were obtained by relaxing the second criterion, meaning that these influencers had fewer than the country's average followers but more than 2,000 followers. For the second batch, 6 influencers (3 from Kenya and 3 from South Africa) were obtained from the sample that took into account all the inclusion criteria, while 23 influencers (2 from Kenya and 21 from South Africa) were obtained by relaxing the third criterion. Additionally, 16 influencers (4 from Kenya and 12 from South Africa) were obtained by relaxing the second criterion. Finally, 65 influencers (49 from Kenya and 16 from South Africa) were recruited by Africa Check through Twitter Ads. Table \ref{tab:reports_tab} summarizes this information. 


Table \ref{tab:desc_stat} provides a summary of the number of influencers and their followers in each country across different stages of our study. The table includes the number of influencers engaged in each phase—pilot, first batch, and second batch—and the total unique followers these influencers had at each stage. Additionally, the table categorizes followers based on the strength of their interactions with the influencers. Followers are classified into three categories: very strong ties, strong ties, and weak ties. Very strong ties are defined as followers who both liked and replied to the influencers' content approximately a year before the intervention began, strong ties are followers who either liked or replied to the content, and weak ties are those who did not interact with the influencers prior to the intervention. 

\begin{table}[!htpb]
    \centering
     \caption{Influencers and Followers Across Study Phases}
    \label{tab:desc_stat}
    \begin{tabular}[t]{lcccccc}
\hline
\textbf{Country} & \textbf{Stage} & \textbf{Influencers} & \textbf{Unique Followers} & \textbf{Very Strong} & \textbf{Strong}\\
\hline
Kenya & Pilot & 20 & 231,874 & 2,021 & 9,188\\
 & First & 76 & 1,310,319 & 19,967 & 81,876 \\
 & Second & 58 & 532,050 & 7,864 & 39,634 \\
\hline
South Africa & Pilot & 20 & 533,812 & 9,323 & 45,559\\
 & First & 76 & 597,357 & 7,955 & 38,176 \\
 & Second & 52 & 620,852 & 4,195 & 23,914 \\
 \hline
\end{tabular}
\end{table}

The influencers have received small payments for the cost of their time and internet data costs. However, the influencers were \textit{not} be contracted based on actually sharing the provided materials. Influencers assigned to the \textit{control} group were asked to avoid sharing content by our institutional partners, meaning that the followers of control influencers were very unlikely to receive the information disseminated by treated influencers through control influencers. 

We additionally seek to compare whether social media micro-influencers are more effective than regular social media ads in ensuring that fact-checking messages reach and are internalized by readers. To assess this, we further cross-randomized some of the followers, irrespective of whether they follow \textit{treated} or \textit{control} social media micro-influencers, into two groups. The first group were \textit{treated} with targeted social media ads that deliver the institutional partners' videos and fact-checks. The ads were disseminated using Twitter's ad targeting tool around the same days that micro-influencers were asked to promote Africa Check's weekly content. The second group are  \textit{control} followers who were not targeted. In sum, the social media influencers' followers have been either  (1) treated with information from the influencers; (2) treated with information from ads; (3) treated with both; or (4) receive no information\footnote{Refer to Appendix \ref{app:foll_rand} for more detail on this.}. 

Importantly, we focused our measurement on, and targeted Twitter ads to, followers who we identify as having very strong or strong ties with at least one social media influencers, and small randomly selected share of followers with weak ties with influencers. The logic of mostly restricting ads to followers with very strong and strong ties reflects the fact that it is impossible to target a meaningful share of followers without a considerable budget, but also that, as pilot data indicates, the Twitter algorithm is much more likely to show the posts of social media influencers to followers with whom they have stronger ties. For this reason, our data collection efforts, and consequently also Twitter ads, mostly focused on users with strong and very strong ties with SMIs.

The study's timing has been as follows. At the start of the study, social media influencers were recruited into the study by fact-checking partners and the research team conducted the randomization. Over the course of eight weeks (and occurring twice a week), \textit{treated} social media influencers received materials from Africa Check and were encouraged to tweet about these materials to their followers. The research team tracked selected followers' social media activity throughout \textcolor{red}{and four weeks after} the study %until Twitter discontinued the Academic access to their API on early May
.  %After the eight weeks of treatment, the research team has approached their followers to take a brief survey. 

Tables \ref{tab:SA_1}-\ref{tab:KE_3} in the Appendix \ref{app:pilot_content} show the content used for the first and second batch intervention in Kenya and South Africa; only the fact-checks were updated in the full intervention launch relative to the pilot content.


 In week 1, the first influencer post or ad is a two-minute 20-second video with %five 
important considerations for people to consider before forwarding a message. %The information source should be clear since news with no information on who wrote the message or where claims come from is likely to be misinformation. Similarly, the information should be verifiable. The information should not elicit emotions such as fear and anger since misinformation is design to prime emotions so that to enhance sharing. Similarly, the use of shocking pictures, videos, or audios is likely to be misleading. Lastly, people should wonder if they think the information is asbolutely true.
This post or ad is accompanied by a second one presenting a fact-check that varies by country. In Kenya, the fact-check debunks the claim that the crime has increased significantly in Nairobi by using a video from Guyana. In South Africa, the fact-check debunks claims about race and private sector ownership.

In week 2, the first influencer post or ad is a 24-second video explaining that the speed at which the COVID-19 vaccine did not compromise its safety. %It highlights that COVID-19 vaccines have been in development for more than 50 years, but until now, there was no funding to create a viable vaccine. Existing research and the work of scientist around the globe was used to develop a safe vaccine. 
This post or ad is accompanied by a second one presenting a fact-check that varies by country. In Kenya, the fact-check debunks the claim that Bill Gates has vowed to pump mRNA vaccines into food supply to 'force-jab' the unvaccinated. In South Africa, the fact-check debunks the claim that South Africans were given a wrong Covid vaccine.

In week 3, the first influencer post or ad is a two-minute 19-second video video explaining how to verify images and videos on social media. %These include reverse image searches on Google or TinEye, looking for visual clues and mapping them, e.g., on wikimapia.org, and double checking that a website's source code has not been manipulated by simply accessing it. 
This post or ad is accompanied by a second one presenting a fact-check using these techniques that varies by country. In Kenya, the fact-check debunks a video showing a girl being assaulted claiming that the incident took place in Kenya. In South Africa, the fact-check debunks the claim that a photo shows the British ambassador to Somalia with a black child in a cage. 

In week 4, the first influencer post or ad is a two-minute 19-second video video explaining how to thoughtfully approach and correct loved ones who share misinformation without upsetting them. %After verifying the misinformation, it is important to first consider and understand their reasons and then communicate privately to avoid embarrassing them. In addition, it is essential to try to avoid conflict and share information on how they can verify claims themselves. 
This post or ad is accompanied by a second one presenting a fact-check using these techniques that varies by country. In Kenya, the fact-check debunks the claim that road crashes claimed more lives than Covid-19. In South Africa, the fact-check debunks power utility Eskom being overstaffed by 66\%.

In week 5, the first influencer post or ad is a 31-second video video explaining why people create and disseminate health misinformation, and the harms associated with its dissemination.  
% Nothing we can really test people on easily. 
This post or ad is accompanied by a second one presenting a fact-check that debunks that expensive treatments advertised on social media can cure HIV.

In week 6, the first influencer post or ad is a one-minute 58-second video video explaining tips for spotting false information and stopping its spread.
% In terms of spotting misinformation, Google the story and see if reputable websites are reporting it. Also, see if the website includes a disclaimer that they cannot guarantee that the content is true or that it is harmless satire. Simlarly, check if there is a journalist or author listed. Fake news websites often don't say who wrote the article. If they do, Google the name and see if they are legit. Lastly, type the website address on https://lookup.icann.org/ and see who registered the website, where, and when. Regarding stopping its spread, don't share the news that looks suspicious, share fact checks, and flag fake news on social media.
This post or ad is accompanied by a second one presenting a fact-check using these techniques that varies by country. In Kenya, the fact-check debunks the claim that a Kenyan doctor invented a cure for high blood pressure. In South Africa, the fact-check debunks the claim that South African Minister of Water and Sanitation suggested people shower ``in groups,'’ however dire the country’s water and electricity crises might be. 

In week 7, the first influencer post or ad is a 38-second video video explaining steps verify health misinformation.
% Perform an internet search of the people or organizations making the claim. Similarly, search if reputable news sources are also taking about the claim. Note that personal anecdotes or endorsement do not substitute scientific proof. Search if the claim has been fact-checked.
This post or ad is accompanied by a second one presenting a fact-check debunking the claim that hot pineapple water cures cancer.

Lastly, in week 8, the first influencer post or ad is a one-minute video on how to avoid being fooled by scams and fight them.
%Poor spelling and grammar are good signs of scams. They also tend to abuse capital letters and exclamation marks. They usually ask for money to access some benefit. Job offers and giveaways from legitimate companies don't come with a fee. They use an official name but link to a completely unrelated website. They ask you to like, comment on, and share it. Don't share the things that looks suspicious, and report suspected scams on social media.
This post or ad is accompanied by a second one presenting a fact-check using these techniques that varies by country. In Kenya, the fact-check points to a scam about posts offering soft loans on WhatsApp. In South Africa, the fact-check points to a scam about rings sold with the promise of weight loss or health benefits. 

All influencer posts include the hashtag \#FactsMatter. Following the findings from \cite{bowlesetal2022}, they use empathetic language that highlights that it is understandable that fears led people to fall for and share misinformation. Moreover, they prime the values and social image, e.g., by highlighting the importance of fact-checking before sharing to protect friends and family.

The analysis and results presented in the following sections focus exclusively on the first and second batch phases for two primary reasons. First, the pilot phase was a preliminary trial designed to evaluate the experiment's functionality, identify best practices for success, and understand how to effectively collaborate with our partner institution, Africa Check. Second, the loss of Academic access to the Twitter API significantly complicated data collection, making it particularly challenging to recover pilot phase information through traditional web scraping techniques.